\documentclass[a4paper]{article}

\usepackage{amsfonts}
\usepackage{amsmath}
\usepackage{amssymb}
\usepackage{graphicx}

\begin{document}

\noindent \large Auteur: Sadia Nazary \\
\noindent \large Studierichting: Informatica\\
\noindent \large Oefeningenreeks 6C nr. 4.2\\

\medskip

\normalsize

$\left(x_n\right)_n=
(0.1,1.1,2.1,0.01,1.01,2.01,0.001,1.001,2.001,\ldots)$\\

We bekijken de drie deelrijen:\\

$x_{3k+1}=10^{-k}$\\

$x_{3k+2}=1+10^{-k}$\\

$x_{3k+3}=2+10^{-k}$\\

Dus:\\

$\lim\limits_{k\to+\infty}x_{3k+1}=0$\\

$\lim\limits_{k\to+\infty}x_{3k+2}=1$\\

$\lim\limits_{k\to+\infty}x_{3k+3}=2$\\

Daarom bestaat
$\lim\limits_{n\to+\infty}x_n$
niet.\\

$\operatorname*{adh}\limits_{n\to+\infty}x_n=\{0,1,2\}$\\

$\liminf\limits_{n\to+\infty}x_n=0$\\

$\limsup\limits_{n\to+\infty}x_n=2$\\

\[
\boxed{
\begin{array}{c}
\lim\limits_{n\to+\infty}x_n
\text{ bestaat niet}\\
\operatorname*{adh}\limits_{n\to+\infty}x_n=\{0,1,2\}\\
\liminf\limits_{n\to+\infty}x_n=0\\
\limsup\limits_{n\to+\infty}x_n=2
\end{array}
}
\]

\begin{thebibliography}{999}
%\bibitem{a} Referentie
\end{thebibliography}

\end{document}